\section{Resultados experimentales y aplicaciones}

\subsection{Rendimiento en modelado de lenguaje}

\subsubsection{Conjunto de datos LM1B}

En el conjunto de datos LM1B (longitud de contexto 128), MDLM supera significativamente a todas las líneas base de modelos de difusión discretos previos (incluyendo la línea base más cercana SEDD) y \textbf{acercó la perplexidad de los modelos AR}.

\begin{knowledgebox}{¿Por qué es importante la perplexidad}
La perplexidad es el indicador de evaluación más central en el modelado de lenguaje. La experiencia demuestra que si un modelo logra buenos puntajes de perplexidad en los conjuntos de entrenamiento/validación, su rendimiento en tareas aguas abajo (como programación o razonamiento) también será mejor. Por lo tanto, la perplexidad es una métrica proxy confiable para medir la calidad de los modelos de lenguaje.
\end{knowledgebox}

Cuando la longitud del contexto se aumenta a 1024, MDLM sigue superando a SEDD, pero existe cierta brecha con respecto a los modelos AR. No obstante, dado que las líneas base de modelos de difusión tempranos mostraban una diferencia enorme frente a las líneas base AR, este resultado ya es muy alentador.

\subsubsection{Capacidad de transferencia sin ejemplos}

Después de entrenar el modelo en OpenWebText, se evalúa la verosimilitud sin ejemplos (zero-shot) en 7 conjuntos de datos no vistos: PTB, arXiv, PubMed, entre otros:
\begin{itemize}
    \item MDLM supera consistentemente a SEDD en todos los conjuntos de datos.
    \item En 3 de los 7 conjuntos de datos, el rendimiento de MDLM es \textbf{superior a la línea base AR}.
\end{itemize}

\subsection{Validación a gran escala: LaDa (MDLM de 8B)}

LaDa (Large Language Diffusion Models) es la validación de MDLM a una escala de 8 mil millones de parámetros. Es esencialmente un modelo MDLM de 8B.

\begin{importantbox}{Rendimiento de MDLM a gran escala}
\begin{itemize}
    \item \textbf{Comprensión multidepósito del lenguaje} (MMLU): MDLM \textbf{supera consistentemente la línea base AR} durante todo el proceso de entrenamiento.
    \item \textbf{Tareas de razonamiento}: La curva de expansión del rendimiento de MDLM muestra un crecimiento \textbf{exponencial}, mientras que la línea base AR solo muestra un crecimiento lineal.
    \item \textbf{Preguntas de sentido común y programación}: MDLM es \textbf{emparejado} o \textbf{ligeramente inferior} a la línea base AR.
\end{itemize}
A pesar de que aún existe una brecha en términos de perplexidad, en tareas aguas abajo prácticas, MDLM supera al modelo AR en múltiples aspectos.
\end{importantbox}

\subsection{Más allá del lenguaje: modelado de datos discretos}

\begin{knowledgebox}{Sesgo izquierda-derecha y estructura de datos}
El lenguaje natural tiene un sesgo inherente de izquierda a derecha—esta es la forma en que los humanos hablan y escriben. Por lo tanto, los modelos AR tienen una ventaja natural en tareas de lenguaje. Sin embargo, en áreas como la generación de proteínas, el descubrimiento de fármacos o la generación molecular, los datos \textbf{no presentan sesgo izquierda-derecha}, lo cual es precisamente el escenario donde MDLM destaca.
\end{knowledgebox}

En tareas de generación molecular (artículo GenMol), los modelos basados en MDLM mostraron una ventaja significativa: mediante la regulación de la temperatura de muestreo, se puede obtener una frontera de Pareto entre la \textbf{calidad} y la \textbf{diversidad} de las moléculas. En contraste, la línea base AR solo puede proporcionar un punto fijo de calidad-diversidad e incapaz de realizar este tipo de compensación flexible.

\subsection{Generación controlada}

MDLM es naturalmente adecuado para la generación controlada. La razón es:

\begin{enumerate}
    \item Durante el proceso de eliminación de ruido, el modelo rellena todas las posiciones enmascaradas, produciendo una \textbf{predicción completa} de la muestra final.
    \item De esto se puede obtener información de alto nivel sobre la muestra final.
    \item Utilizando esta información, se pueden guiar muestras con atributos específicos durante el proceso de muestreo.
\end{enumerate}

\begin{warningbox}{Los modelos AR tienen dificultades para realizar muestreo controlado}
Los modelos AR tienen dificultades conocidas en términos de generación controlada. Debido a que, durante el proceso de generación palabra por palabra, el modelo no puede prever la estructura global futura, es muy difícil imponer restricciones globales durante el proceso de generación. El mecanismo de atención bidireccional y la predicción global de eliminación de ruido de MDLM le confieren una ventaja esencial en la generación controlada.
\end{warningbox}

\subsection{Resumen de la sección}

Los resultados experimentales indican: (1) MDLM se acerca al rendimiento de AR a pequeña escala, e incluso supera a AR en tareas aguas abajo a gran escala (especialmente razonamiento); (2) En datos discretos sin sesgo izquierda-derecha (como generación molecular), el rendimiento de MDLM es particularmente destacado; (3) La flexibilidad de muestreo de MDLM lo hace adecuado para aplicaciones como la generación controlada.

\newpage

````
